\section{Background}
\label{sec:background}

This section fixes the terminology used throughout and states the properties of
code property graphs and graph neural networks that the measurements depend on.
Readers familiar with both may proceed to Section~\ref{sec:method}.

\subsection{Code property graphs}
A code property graph~\citep{6956589} represents a function as a single graph
that superimposes three classical program representations. Nodes correspond to
syntactic elements: a call, an identifier, a control structure, a literal. Edges
carry a relation type, and four occur in the exports used here. \emph{AST} edges
connect a syntactic construct to its constituents, so a call node is joined to
its arguments. \emph{CFG} edges connect a statement to those that may execute
next. \emph{DDG} edges connect a definition of a value to the uses that may
observe it. \emph{CDG} edges connect a predicate to the statements whose
execution it governs. The latter two derive from the program dependence
graph~\citep{10.1145/24039.24041}, which predates their use as learning
features by three decades.

Two properties of this representation matter for what follows. First, the edges
are \emph{directed}, and the direction is fixed by the relation: an AST edge runs
from parent to child, a CFG edge from a statement to its successor, a DDG edge
from a definition to a use. A tool consuming the export unaltered therefore
inherits that orientation. Second, the relation types are \emph{labels rather
than separate graphs}: all four occupy one edge set, and distinguishing them
requires a model that can read the label. Section~\ref{sec:graph-worth} measures
what each of these two properties is worth, and finds that they differ sharply:
the orientation is consequential and the labels are not.

The graphs in this work are produced by Joern~\citep{joern_cpg_docs} in its
per-method representation, one graph per function; the representation and its
query model are described further
in~\citep{yamaguchi2020_querying_whole_program_graphs}.

\subsection{Graph neural networks}
A graph neural network computes a representation for each node by repeatedly
aggregating from that node's neighbours. Writing $h_i^{(l)}$ for the
representation of node $i$ at layer $l$, one round of message passing takes the
form
\[
  h_i^{(l+1)} = \phi\Bigl(h_i^{(l)},\; \bigoplus_{j \in \mathcal{N}(i)}
                    \psi\bigl(h_i^{(l)}, h_j^{(l)}, e_{ji}\bigr)\Bigr),
\]
where $\mathcal{N}(i)$ is the set of neighbours from which $i$ may aggregate,
$e_{ji}$ is the feature of the edge joining them, $\bigoplus$ is a
permutation-invariant aggregator, and $\phi, \psi$ are learned. After $L$ layers
a node's representation depends on nodes at most $L$ hops away, and a
permutation-invariant readout pools the node representations into one vector for
classification.

Three consequences are used later. Because $\mathcal{N}(i)$ follows edge
direction, a node aggregates only from its predecessors unless the reverse edges
are supplied explicitly. Because the aggregator is permutation-invariant and
depth-limited, information from a node more than $L$ hops away cannot reach $i$
at all. Raising $L$ is not a free remedy: node representations converge with
depth~\citep{li2018deeper,oono2020graph}, and signal routed through few
intermediate nodes is compressed as the receptive field
widens~\citep{alon2021bottleneck,topping2022understanding}. And because $e_{ji}$ enters $\psi$, an operator that omits the edge term
cannot distinguish relation types: we verify in
Section~\ref{sec:graph-worth} that one of the four architectures studied here
produces bit-identical output when the edge labels of a fixed graph are
permuted. Reading the label requires either an edge-conditioned message or a
per-relation weight matrix, as in relational graph
convolution~\citep{schlichtkrull2018modeling}.

The four aggregation families compared in this paper are attention
(GAT~\citep{DBLP:conf/iclr/VelickovicCCRLB18}), spectral convolution
(GCN~\citep{Kipf2016SemiSupervisedCW}), sum aggregation
(GIN~\citep{xu2018how,morris2019weisfeiler}), and global attention via a graph
transformer~\citep{3454287.3455360}.

\subsection{Evaluating a detector}
A detector is trained on a set of functions labelled vulnerable or not, and
scored on a disjoint set. Every metric reported here is computed on functions
the model has not been trained on; what varies between protocols, and what
Section~\ref{sec:ladder} measures, is what else those functions may share with
the training set. We report ROC-AUC as the primary metric because it is
invariant to class prevalence, which differs across the protocols compared, and
Matthews correlation alongside it because it is not. Comparisons between graph
models are sensitive to exactly these choices, and reported gains have been
shown to shrink or vanish once model selection and splitting are held
fixed~\citep{errica2020fair}; we therefore hold both fixed across every
condition in Section~\ref{sec:results}.
